\documentclass[11pt]{article}
\usepackage[margin=0.93in]{geometry}
\usepackage{amsmath,amssymb,amsthm,mathtools,booktabs,microtype}
\usepackage[colorlinks=true,linkcolor=blue!40!black,urlcolor=blue!50!black,citecolor=blue!40!black]{hyperref}
\usepackage{xcolor}
\newtheorem{theorem}{Theorem}[section]
\newtheorem{lemma}[theorem]{Lemma}
\newtheorem{corollary}[theorem]{Corollary}
\theoremstyle{remark}\newtheorem{remark}[theorem]{Remark}
\newcommand{\Rea}{\operatorname{Re}}
\newcommand{\Ima}{\operatorname{Im}}
\newcommand{\e}{\mathrm{e}}
\newcommand{\ii}{\mathrm{i}}
\newcommand{\G}{\mathcal G}
\newcommand{\GRH}{\mathrm{GRH}}
\newcommand{\ord}{\operatorname{ord}}
\allowdisplaybreaks[2]
\title{\textbf{Peak rigidity and simultaneous recurrence of\protect\\ smoothed prime-counting errors}\\[4pt]\large An extension of the one-sided prime-sum method for Dirichlet GRH}
\author{Research note}
\date{26 November 2026}
\begin{document}
\maketitle
\begin{abstract}
A smoothing kernel used in the Lean proof of Dirichlet GRH has a special property: on the critical line its zero coefficients are positive real numbers. Exploiting this property, we compute the exact asymptotic upper envelope for the normalized error in the identity residue class, and prove quantitative rigidity of the entire vector of errors. Near-maximal identity-class errors force uniform recurrence of the whole limiting error process, including all multiplicative shifts in the argument. The symmetrized errors satisfy a stronger linear deficit bound and an associated positive-semidefinite matrix inequality. The results hold simultaneously for every fixed finite collection of moduli, and a single sequence of integers realizes the canonical peak profiles for all moduli and all periodic arithmetic weights. No assumption of linear independence or simple zeros is made. We also prove a converse one-sided criterion by Landau's method. The only nonclassical input is Dirichlet GRH, the endpoint of the supplied Lean project; its proof is not independently replayed here.
\end{abstract}

\section{Source analysis and the research question}
The supplied repository \texttt{lsernau/grh-lean} \cite{repo} is a compact and a readable rewriting of Levent Alp\"oge's proof. Its README describes Lean 4.35.0, Mathlib commit \texttt{56a00585} (full hash in \texttt{SOURCES.md}), five endpoint theorems, and checks with the Lean kernel and Comparator. Its authors report only the standard axioms \texttt{propext}, \texttt{Classical.choice}, and \texttt{Quot.sound}. The contemporaneous Wikipedia article \cite{wiki} records the announced proof and independent machine verification, while noting that human exposition remains incomplete. The present argument uses the declared mathematical endpoint, not a new verification of the repository.

The compact project's \texttt{src/New/Main.lean} assembles three geometric and arithmetic bounds on theta lifts into \texttt{star\_inequality}, then deduces \texttt{eventually\_primeError\_le\_rpow} and \texttt{New.grh}. In the readable project, \texttt{src/New/Reduction/Zeros/Criterion.lean} states a one-sided nonvanishing criterion with the weighted remainder
\[
 \sum_{\substack{n\le x\\n\equiv1\ (q)}}\Lambda(n)(1-\chi_d(n))
       \left(\frac1n-\frac1x\right)-\frac{\log x}{\varphi(q)}-c,
\]
where $\chi_d$ is a quadratic character. The pole computation in that file features $1/[s(1-s)]$. The compact file \texttt{src/Lib/Analytic/Landau/PolarLandau.lean} treats one-sided bounds on Laplace transforms with nonnegative real residues. The Wikipedia discussion of explicit formulas and prime-distribution bounds supplies the other motivation.

\emph{Question.} Once all the zeros lie on the critical line, can this particular smoothing kernel yield not only bounds but an exact joint extremal law for the prime errors? The answer below is yes. The key observation is
\begin{equation}\label{eq:positive}
 \frac1{\rho(1-\rho)}=\frac1{1/4+\gamma^2}>0
 \qquad (\rho=1/2+\ii\gamma).
\end{equation}
Unlike a usual explicit formula for unsmoothed prime counts, all spectral contributions in the identity class can align at a common phase. The precise new target is the uniform translate inequality in Theorem~\ref{thm:main}, its inverse-class refinement, and the universal periodic-weight recurrence. The one-sided criterion itself is a classical Landau-type argument, not the main novelty claim.

\section{Notation and exact envelope}
Fix $q\ge1$, let $G=(\mathbb Z/q\mathbb Z)^\times$ with $h=\varphi(q)$ (for $q=1$ take the trivial group), and write $\widehat G$ for its Dirichlet characters extended by zero on nonunits. Put $\delta_\chi=1$ for the principal character and $0$ otherwise. Define
\begin{align}
 S_\chi(x)&=\sum_{n\le x}\Lambda(n)\chi(n)(1/n-1/x),\nonumber\\
 c_\chi&=\lim_{s\to1}\left(-\frac{L'}{L}(s,\chi)-\frac{\delta_\chi}{s-1}\right)-\delta_\chi,\label{eq:constant}\\
 E_\chi(x)&=\sqrt{x}\big(S_\chi(x)-\delta_\chi\log x-c_\chi\big).
\end{align}
The limit is defined by analytic continuation; $L(1,\chi)\ne0$ for nonprincipal characters. For the principal character $c_{\chi_0}=-1-\gamma_E-\sum_{p\mid q}\log p/(p-1)$; here $\gamma_E$ is Euler's constant. For $a\in G$, set
\begin{align}
 c_a&=h^{-1}\sum_\chi\overline{\chi(a)}c_\chi,\\
 E_a(x)&=\sqrt{x}\left[\sum_{\substack{n\le x\\n\equiv a\ (q)}}\Lambda(n)(1/n-1/x)
                      -h^{-1}\log x-c_a\right]\nonumber\\
 &=h^{-1}\sum_\chi\overline{\chi(a)} E_\chi(x).\label{eq:fourier}
\end{align}
These are real for residue classes. Subscripts $q$ will be included when comparing different moduli. Let $\chi^*$ be the primitive character inducing $\chi$, with conductor $f_\chi$. In sums over zeros, use the nontrivial zeros of $L(s,\chi^*)$, counted with multiplicity; finite Euler-factor zeros on $\Rea s=0$ are not included.

Under $\GRH$ for the characters modulo $q$, all these zeros have the form $1/2+\ii\gamma$. Define the finite positive mass, Fourier series, and peak vector
\begin{align}
 M_\chi&=\sum_{\gamma\ (\chi^*)} \frac1{1/4+\gamma^2},&
 B_\chi(t)&=\sum_{\gamma\ (\chi^*)}\frac{\e^{\ii\gamma t}}{1/4+\gamma^2},\label{eq:spectrum}\\
 B_a(t)&=h^{-1}\sum_\chi\overline{\chi(a)}B_\chi(t),&
 v_a&=B_a(0),\nonumber\\
 D(t)&=v_1-B_1(t),& r(x)&=(1+\log x)/\sqrt{x}.\label{eq:deficit}
\end{align}
The standard zero-counting estimate $N_\chi(T)=O_\chi((T+1)\log(T+2))$ for $|\gamma|\le T$ implies absolute uniform convergence of \eqref{eq:spectrum}. Every $M_\chi>0$ since primitive Dirichlet $L$-functions have infinitely many nontrivial zeros. Conjugation yields $B_{\bar\chi}(t)=\overline{B_\chi(t)}$, $M_{\bar\chi}=M_\chi$, so $B_a$ and $v_a$ are real, and $v_a=v_{a^{-1}}$. In particular $v_1=h^{-1}\sum M_\chi>0$.

\begin{theorem}[Peak rigidity]\label{thm:main}
Assume $\GRH$ for all characters modulo $q$. Then:
\begin{enumerate}
\item For $x\ge2$,
\begin{equation}\label{eq:approx}
 E_a(x)=B_a(\log x)+O_q(r(x)),\qquad
 \limsup_{x\to\infty} E_1(x)=v_1.
\end{equation}
Moreover $\limsup_{x\to\infty}|E_a(x)|\le v_1$ for every $a$.
\item The scalar deficit $D(t)$ is nonnegative for all real $t$, and for all $a$ and all real $t,u$,
\begin{equation}\label{eq:translate}
 \boxed{\ |B_a(t+u)-B_a(u)|^2\le 2v_1 D(t).\ }
\end{equation}
In particular the entire limiting error process approximately repeats whenever $E_1$ nearly attains its sharp upper envelope.
\item The sharper inverse-class and Fourier-weighted bounds are
\begin{align}
 \left|\frac{B_a(t)+B_{a^{-1}}(t)}2-v_a\right|&\le D(t),\label{eq:linear}\\
 \frac1h\sum_\chi\frac{|B_\chi(t)-M_\chi|^2}{M_\chi}&\le 2D(t).\label{eq:ellipsoid}
\end{align}
For the real symmetric matrices indexed by $a,b\in G$,
\begin{align*}
 V_{a,b}&=v_{ab^{-1}},& H(t)_{a,b}&=\frac{B_{ab^{-1}}(t)+B_{ba^{-1}}(t)}2,
\end{align*}
$V$ is positive definite, $V-H(t)$ is positive semidefinite, and its trace is $hD(t)$.
\item For every $\epsilon>0$, the set of $t$ for which
\[
 \sup_{a\in G,\ u\in\mathbb R}|B_a(t+u)-B_a(u)|<\epsilon
\]
is relatively dense (every interval of some fixed length meets it). Consequently the near-peak scales of $E_1$ are relatively dense on the logarithmic axis after ignoring a bounded initial interval.
\end{enumerate}
\end{theorem}

\begin{corollary}[Unique exposed peak and local universality]\label{cor:local}
For any sequence $x_j\to\infty$, the following are equivalent:
\begin{enumerate}
\item $E_1(x_j)\to v_1$;
\item $(E_a(x_j))_{a\in G}\to(v_a)_{a\in G}$;
\item $E_a(x_j\e^u)\to B_a(u)$ for every $a$, uniformly for $u\ge-\tfrac12\log x_j$.
\end{enumerate}
In particular the hyperplane with identity coordinate $v_1$ meets the asymptotic cluster set of $(E_a(x))$ at a single point. The same claims about limiting peaks apply if $x$ is restricted to positive integers.
\end{corollary}

\section{Proof of the spectral expansion}
We give the normalization and sign explicitly because using the wrong sign would exchange the upper and lower extrema. For a primitive character $\eta$ of conductor $f$, including the zeta character of conductor $1$, write $A(s)=-L'(s,\eta)/L(s,\eta)$. For $c>1$, Mellin inversion gives
\begin{equation}\label{eq:perron}
 S_\eta(x)=\frac1{2\pi\ii}\int_{c-\ii\infty}^{c+\ii\infty}
      A(s)\frac{x^{s-1}}{s(s-1)}\,ds.
\end{equation}
Indeed $A(s)=\sum_{n\ge1}\Lambda(n)\eta(n)n^{-s}$ absolutely there, and the integral of $y^s/[s(s-1)]$ equals $y-1$ if $y>1$ and $0$ if $0<y\le1$. This follows by moving the line left across $0$ and $1$, or right across no poles. At $y=1$ the kernel is continuous, so no half-weight convention is needed. Multiplying by $1/x$ and setting $y=x/n$ proves \eqref{eq:perron}; the interchange is justified by the integrable $O((1+|\Ima s|)^{-2})$ kernel.

Move the contour to $\Rea s=-1/2$. We use standard analytic facts for primitive Dirichlet $L$-functions (see Davenport \cite{davenport}, Montgomery--Vaughan \cite{mv}): the functional equation, zero-free line $\Rea s=1$, simple pole of $\zeta$ at $1$, the zero-counting bound above, and the local log-derivative estimate on bounded vertical strips at large height
\[
 \frac{L'}{L}(\sigma+\ii T,\eta)
 =\sum_{|\Ima\rho-T|\le1}\frac{1}{\sigma+\ii T-\rho}
   +O_\eta(\log(|T|+2)),\quad -1/2\le\sigma\le c.
\]
The local number of zeros is $O_\eta(\log(|T|+2))$. Thus one can choose heights separately tending to $+\infty$ and $-\infty$ with distance at least $c_\eta/\log(|T|+2)$ from every nearby ordinate: remove intervals of this radius around the $O_\eta(\log |T|)$ ordinates meeting a unit interval, choosing $c_\eta$ small. On the horizontal edges the log derivative is $O_\eta(\log^2 |T|)$ and their contributions tend to zero for each fixed $x$. The functional equation and the absolutely convergent log derivative on $\Rea s=3/2$ give $A(-1/2+\ii t)=O_\eta(\log(|t|+2))$ (Stirling for the gamma factors); there is no zero on the new line. The new vertical integral is absolutely convergent and $O_\eta(x^{-3/2})$ uniformly for $x\ge2$.

The residues crossed are:
\begin{itemize}
\item At $s=1$, the kernel supplies a simple pole. If $\eta$ is nonprincipal its contribution is $A(1)$. For the zeta character, $A(s)=1/(s-1)-\gamma_E+O(s-1)$; multiplication by $x^{s-1}/[s(s-1)]$ gives residue $\log x-1-\gamma_E$. Together these are $\delta_\eta\log x+c_\eta$.
\item At each nontrivial zero $\rho$ of multiplicity $m$, $A(s)$ has residue $-m$. The integrand residue is $m x^{\rho-1}/[\rho(1-\rho)]$.
\item At $s=0$, $A$ has at most a simple pole, so the residue is $x^{-1}$ times a polynomial in $\log x$ of degree at most one. There are no other trivial zeros in the strip crossed.
\end{itemize}
The zero sum converges absolutely even before GRH for each fixed $x$, since the real parts of nontrivial zeros are in $(0,1)$ and the zero-counting bound controls the $O(|\gamma|^{-2})$ tail. Under GRH the residue identity becomes
\[
 S_\eta(x)=\delta_\eta\log x+c_\eta
   +x^{-1/2} B_\eta(\log x)+O_\eta((1+\log x)/x).
\]
This explicitly treats central zeros and arbitrary multiplicities.

For an imprimitive $\chi$ induced by $\chi^*$, the only difference between the prime-power sums is the removal of the primes $p\mid q$, $p\nmid f_\chi$. For each such prime the removed contribution is
\[
 \sum_{p^k\le x}\chi^*(p)^k\log p(1/p^k-1/x)
 =\frac{\chi^*(p)\log p}{p-\chi^*(p)}+O_p((1+\log x)/x).
\]
The bound follows from the geometric tail and $\lfloor\log x/\log p\rfloor$ terms of the $1/x$ sum. The constant equals the correction in \eqref{eq:constant} by differentiating the removed Euler factor. Therefore $E_\chi(x)=B_\chi(\log x)+O_q(r(x))$; character orthogonality proves the first equality of \eqref{eq:approx}.

\section{Proof of peak rigidity and recurrence}
Write $w_\gamma=(1/4+\gamma^2)^{-1}$, retaining multiplicities. The following identity, with no asymptotic error, is the heart of the argument:
\begin{equation}\label{eq:energy}
 D(t)=\frac1h\sum_\chi\sum_\gamma w_\gamma(1-\cos(\gamma t))
 =\frac1{2h}\sum_\chi\sum_\gamma w_\gamma|\e^{\ii\gamma t}-1|^2.
\end{equation}
Here $B_1$ is real by conjugation. In particular $D(t)\ge0$ and $|B_a(t)|\le v_1$. For every $a,u$, write the translate difference as
\[
 \frac1h\sum_{\chi,\gamma}w_\gamma\overline{\chi(a)}\e^{\ii\gamma u}(\e^{\ii\gamma t}-1).
\]
The other phase has absolute value $1$. Cauchy--Schwarz for the finite measure assigning mass $w_\gamma/h$ gives \eqref{eq:translate}, since the total mass is $v_1$ and the second moment is $2D(t)$. Absolute convergence makes all rearrangements legitimate.

For each character set $D_\chi(t)=M_\chi-\Rea B_\chi(t)\ge0$. Their average is $D(t)$. Also
\[
 |M_\chi-B_\chi(t)|^2
 \le M_\chi\sum_\gamma w_\gamma|1-\e^{\ii\gamma t}|^2=2M_\chi D_\chi(t).
\]
Summing yields \eqref{eq:ellipsoid}. In the symmetrization, conjugate-character pairing gives
\[
 v_a-\tfrac12(B_a(t)+B_{a^{-1}}(t))
 =h^{-1}\sum_\chi \Rea\chi(a)D_\chi(t).
\]
Since $|\Rea\chi(a)|\le1$, this proves \eqref{eq:linear}. The vectors $(\overline{\chi(a)})_{a\in G}$ are a basis of eigenvectors for the displayed convolution matrices. Their eigenvalues for $V$ are $M_\chi$ and for $H$ are $\Rea B_\chi(t)$, respectively. This proves the definiteness assertions and trace identity. In particular, the positive cone of deficit matrices is a simultaneous constraint across residue classes, not just a collection of separate estimates.

\subsection*{Recurrence without independent zeros}
For finitely many frequencies $\gamma_1,\ldots,\gamma_k$, let $g(t)=(\e^{\ii\gamma_1t},\ldots,\e^{\ii\gamma_kt})$. Its orbit closure in the torus is a compact subgroup $K$. For any symmetric open neighborhood $U$ of the identity in $K$, translates $g(s)U$ for $s\in\mathbb R$ cover $K$ by density. A finite subcover $g(s_1)U,\ldots,g(s_m)U$ exists. For any $t$, some $g(t-s_j)$ belongs to $U$. If $A\le s_j\le B$ for all $j$, there is such a return in $[t-B,t-A]$; hence returns are relatively dense. If $A=B$, enlarge the interval length to $1$. No density in the full torus is required.

For the full family, discard a finite tail of arbitrarily small mass in \eqref{eq:energy}. The discarded tail contributes at most twice its mass to $D(t)$. Return sufficiently close to identity on the retained frequencies to make the finite part equally small. This makes $D(t)$ arbitrarily small on a relatively dense set, and \eqref{eq:translate} proves the stated uniform recurrence. As such returns occur arbitrarily far to the right, the upper bound $B_1(t)\le v_1$ is attained as a limsup. With \eqref{eq:approx} this completes Theorem~\ref{thm:main}.

For Corollary~\ref{cor:local}, if $E_1(x_j)\to v_1$, then $D(\log x_j)\to0$. Equation \eqref{eq:translate} gives a uniform bound in all $u$; \eqref{eq:approx} adds an error uniformly tending to zero for $x_j\e^u\ge\sqrt{x_j}$. The other implications follow by taking $u=0$ and $a=1$. Since all $B_a$ are uniform limits of continuous finite exponential sums, they are uniformly continuous. Rounding $\e^t$ to its nearest positive integer changes its logarithm by $o(1)$; the errors in \eqref{eq:approx} also tend to zero. This proves the integer statement.

\section{Arithmetic evaluation of the peak vector}
The masses require no enumeration of zeros. For a nonprincipal primitive $\eta$ of conductor $f$, let $b\in\{0,1\}$ be its parity ($\eta(-1)=(-1)^b$), and let
\[
 \Xi_\eta(s)=(f/\pi)^{(s+b)/2}\Gamma((s+b)/2)L(s,\eta).
\]
This is entire of order one with zeros precisely the nontrivial zeros. For the zeta character take $\Xi_1(s)=\xi(s)=\tfrac12 s(s-1)\pi^{-s/2}\Gamma(s/2)\zeta(s)$. The genus-one Hadamard product and logarithmic differentiation give
\begin{equation}\label{eq:mass}
 \frac{\Xi_\eta'}{\Xi_\eta}(1)-\frac{\Xi_\eta'}{\Xi_\eta}(0)
 =\sum_\rho\left(\frac1{1-\rho}+\frac1\rho\right)=M_\eta.
\end{equation}
The left difference cancels the constant in the logarithmic derivative; the right sum is absolutely convergent. Under GRH the summands are positive by \eqref{eq:positive}. The functional equation $\Xi_\eta(s)=\varepsilon_\eta\Xi_{\bar\eta}(1-s)$ and conjugation imply that the difference is $2\Rea(\Xi_\eta'/\Xi_\eta)(1)$. Thus, with $\psi=\Gamma'/\Gamma$,
\begin{equation}\label{eq:massformula}
 M_\chi=\begin{cases}
 2+\gamma_E-\log(4\pi),& f_\chi=1,\\
 \log(f_\chi/\pi)+\psi((1+b)/2)
 +2\Rea\dfrac{L'}{L}(1,\chi^*),& f_\chi>1.
 \end{cases}
\end{equation}
For the first line use $\xi'/\xi(1)=1+\gamma_E-\tfrac12\log\pi+\tfrac12\psi(1/2)$ and $\psi(1/2)=-\gamma_E-2\log2$. Imprimitive characters use the mass of their primitive inducing character; their centering constants still involve the removed factors. The exact envelope $v_1$ and every $v_a$ are obtained from \eqref{eq:massformula} by a finite inverse Fourier transform. As $M_\chi>0$ for every character, the triangle inequality is strict in $|v_a|<v_1$ for $a\ne1$ (the trivial character has phase $1$, and some other character separates $a$ from $1$).

\begin{table}[ht]
\centering
\caption{Approximate peak vectors; residue classes in increasing order.}\label{tab:examples}
\begin{tabular}{rll}
\toprule
$q$ & $a$ & $v_a$\\\midrule
1 & identity & $0.046191417932$\\
3 & 1, 2 & $0.079710693895,\ -0.033519275963$\\
5 & 1, 2, 3, 4 & $0.152297809271,\ -0.027591384016,$\\
  & & $-0.027591384016,\ -0.050923623308$\\
\bottomrule
\end{tabular}
\end{table}
For illustration, the associated program evaluates the constants by Hurwitz zeta Laurent coefficients, independently of zero lists. For a nonprincipal character modulo a prime $q$, one has
\[
 L(1,\chi)=-q^{-1}\sum_{a=1}^{q-1}\chi(a)\psi(a/q),\quad
 L'(1,\chi)=-q^{-1}\sum_{a=1}^{q-1}\chi(a)\gamma_1(a/q)-\log q\,L(1,\chi),
\]
where $\gamma_1(z)$ is the first generalized Stieltjes constant. These follow from the Laurent series of Hurwitz zeta at $1$ and $\sum\chi(a)=0$. The accompanying CSV contains additional moduli and smoothed prime-power samples. The decimal numbers illustrate the formulas; they are not rigorous interval enclosures. In particular, at finite $x$ the term $O_q(r(x))$ may let an observed identity-class error exceed $v_1$; the upper envelope is asymptotic.

\section{One sequence for all moduli and periodic weights}
\begin{theorem}[Universal simultaneous peaks]\label{thm:univ}
Assume Dirichlet GRH. There is an increasing sequence of integers $n_j\to\infty$ such that for every fixed $q\ge1$ and $a\in(\mathbb Z/q\mathbb Z)^\times$,
\[
 E_{q,a}(n_j\e^u)-B_{q,a}(u)\longrightarrow0
 \qquad\text{uniformly for }u\ge-\tfrac12\log n_j.
\]
For any finite collection of moduli and fixed tolerance, the corresponding approximate periods of their limiting processes are relatively dense on the logarithmic axis.
\end{theorem}
\begin{proof}
Take a finite union of all retained zero frequencies for the finitely many moduli in the compact-group recurrence argument. The argument applies to each of their tails and total masses simultaneously. For each $j$, choose a sufficiently large return $t_j$ for the moduli $1,\ldots,j$ so that their uniform translate errors are $<1/(2j)$. Enlarge $t_j$ to make the arithmetic error \eqref{eq:approx} uniformly $<1/j$ for all those moduli on $x\ge\tfrac12\e^{t_j/2}$, and to make rounding $\e^{t_j}$ to $n_j$ perturb all the finitely many limiting series by $<1/(2j)$ (uniform continuity). The return set is unbounded so all these requirements can be imposed, including increasing $n_j$. For each fixed $q$ and sufficiently large $j$, the bound is $O(1/j)$ on the stated range. This establishes the theorem without any independence hypothesis, and without a uniform-in-$q$ rate.
\end{proof}

\begin{corollary}[Arbitrary periodic weights, including quadratic twists]\label{cor:periodic}
Let $F:\mathbb Z\to\mathbb C$ be periodic of period $Q\ge1$. Define
\begin{align*}
 \mu_F&=\varphi(Q)^{-1}\sum_{a\in G_Q}F(a),\\
 c_F&=\sum_{a\in G_Q}F(a)c_{Q,a}
      +\sum_{p\mid Q}\log p\sum_{k\ge1}F(p^k)/p^k,\\
 E_F(x)&=\sqrt{x}\left(\sum_{n\le x}\Lambda(n)F(n)(1/n-1/x)-\mu_F\log x-c_F\right).
\end{align*}
Then $E_F(x)=B_F(\log x)+O_{F,Q}(r(x))$, where $B_F(u)=\sum_{a\in G_Q}F(a)B_{Q,a}(u)$. Along the same integer sequence as Theorem~\ref{thm:univ}, $E_F(n_j\e^u)\to B_F(u)$ uniformly for $u\ge-\tfrac12\log n_j$. At a single master-modulus scale, the spectral translate error for $F$ is bounded by $\sum_{a\in G_Q}|F(a)|\sqrt{2v_{Q,1}D_Q(t)}$.
\end{corollary}
\begin{proof}
The contributions from unit classes follow by linearity. For nonunit $n$ with $\Lambda(n)\ne0$, $n=p^k$ for some $p\mid Q$. The geometric tail and the number of terms give the additional constant and an $O_{F,Q}((1+\log x)/x)$ remainder before normalization. Apply Theorem~\ref{thm:univ} and the triangle inequality to the finitely many unit classes. This includes $F(n)=\mathbf1_{n\equiv1\ (q)}(1-\chi_d(n))$ of period dividing $qd$, the twisted weight from the repository. For such a signed Fourier combination the canonical recurring value need not be its maximum: the positive coefficients needed for the exact upper-envelope assertion may fail. The recurrence assertion still holds.
\end{proof}

\section{A converse via one-sided Landau theory}
To clarify the relation to the repository, we show that a bounded upper envelope in the identity class is already enough to recover GRH for that modulus. This section does not assume GRH.

\begin{theorem}[One-class criterion]\label{thm:converse}
Fix $q\ge1$. Dirichlet GRH for the characters modulo $q$ (equivalently for their primitive inducing characters) holds if and only if
\[
 \limsup_{x\to\infty} E_1(x)<+\infty.
\]
If GRH fails for any of these characters, then the limsup is $+\infty$. If GRH holds, it equals the explicit mass average $v_1$ from \eqref{eq:massformula}.
\end{theorem}
\begin{proof}
The forward implication was proved in Theorem~\ref{thm:main}. Conversely put $R(x)=E_1(x)/\sqrt{x}$. With no GRH assumption, for $\Rea s>1$ termwise integration shows
\begin{align}
 T(s)&=\int_1^\infty R(x)x^{-s}\,dx\nonumber\\
 &=\frac1h\sum_\chi\frac{-L'(s,\chi)/L(s,\chi)}{s(s-1)}
 -\frac{1/h}{(s-1)^2}-\frac{c_1}{s-1}.\label{eq:mellin}
\end{align}
Indeed the integral of $(1/n-1/x)\mathbf1_{x\ge n}x^{-s}$ is $n^{-s}/[s(s-1)]$. The main terms cancel both pole orders at $s=1$ by \eqref{eq:constant}. Equation \eqref{eq:mellin} extends meromorphically to $\Rea s>1/2$, with possible poles only at nontrivial zeros of the characters. At any such $\rho$ the residue is
\[
 \frac{\sum_\chi \ord_\rho L(s,\chi)}{h\rho(1-\rho)}\ne0.
\]
There is no cancellation, because all multiplicities in the numerator are nonnegative, and at real poles in $(1/2,1)$ the residue is strictly positive.

Suppose the limsup is finite (including $-\infty$). Choose a real $C$ with $R(x)\le Cx^{-1/2}$ for every $x\ge1$, by increasing $C$ on a compact initial interval; $R$ is continuous there. The nonnegative function $g(x)=Cx^{-1/2}-R(x)$ has Mellin transform
\[
 U(s)=\int_1^\infty g(x)x^{-s}\,dx=\frac{C}{s-1/2}-T(s)\qquad(\Rea s>1).
\]
Let $\sigma_c$ be the abscissa of convergence of the positive integral, possibly $-\infty$. It is at most $1$ because the elementary bound on $S_\chi$ is $O((1+\log x)^2)$. For a nonnegative function, a finite abscissa of convergence is a singular point of its transform on the real axis (Landau's lemma; see below). If $\sigma_c>1/2$, the continuation of $U$ is either holomorphic there, impossible by the lemma, or has a simple pole with strictly negative residue by \eqref{eq:mellin}. A negative residue would force $U(\sigma)\to-\infty$ as real $\sigma\downarrow\sigma_c$, contradicting the positive integral. Hence $\sigma_c\le1/2$, and the integral is holomorphic in the open half-plane $\Rea s>1/2$. By uniqueness of analytic continuation there can be no pole there, hence no zero in that half-plane. The functional equation for the primitive characters and closure of the character family under conjugation give the critical-line conclusion.\end{proof}

\paragraph{Landau lemma (proof).} Substitute $x=\e^t$, so the integral is the Laplace transform of a nonnegative locally integrable function $e^t g(e^t)$. If its finite abscissa $\sigma_c$ were a regular point, choose a real $s_0>\sigma_c$ so close that the Taylor disk at $s_0$ extends to some $s_0-d<\sigma_c$. This is possible by a disk of holomorphy about $\sigma_c$. Differentiation under the integral for $s_0>\sigma_c$ is valid by exponential decay relative to a smaller convergent real parameter. The Taylor series at $s_0-d$ has nonnegative terms, since the $k$th derivative has sign $(-1)^k$. By monotone convergence the sum of the convergent Taylor series equals the positive integral with parameter $s_0-d$. This contradicts the definition of the abscissa. Holomorphy in the half-plane of convergence follows from dominated convergence on compact subsets. This proves the lemma and exhibits exactly the sign mechanism also exploited by the source project.

\section{Novelty, verification, and further research}
\paragraph{What is and is not new.} The ingredients---Riesz/Perron explicit formulas, the Hadamard mass identity, compact-group recurrence, and Landau's lemma---are classical. Suzuki \cite{suzuki} studies related sign criteria for weighted Mangoldt sums; the original GRH proof uses windowed one-sided bounds, and the supplied rewrite uses twists. Rubinstein--Sarnak \cite{rs} study prime races and their limiting distributions, often with additional hypotheses about zero ordinates. The specific proposed contribution here is the exact simultaneous extremal and uniform-translate control for the $1/n-1/x$ kernel, especially the linear bound on inversion-symmetric errors, the positive deficit matrix, and the transfer to every periodic weight at one sequence of scales. A targeted search found related one-sided criteria and prime-race literature but did not establish priority for this precise combination. It is a proved candidate for a new result, not a claim that a long-standing named conjecture has been resolved or that the result has been peer reviewed.

\paragraph{Verification.} The mathematical proofs are given above; the analytic inputs are standard Dirichlet $L$-function theory and the stated GRH endpoint. The attached Python program verifies the finite phase-energy identity and inequality using exact rational arithmetic and computes sample constants and prime-power sums. The finite tests are not a proof of the infinite analytic statements. No Lean executable was available in the working environment, and the upstream repository was not rebuilt. A potential Lean development would first formalize the character-weighted positive phase-energy inequality for finite sums, then extend via summability, prove the explicit formula and conductor correction, and obtain recurrence by compactness. The only nonclassical analytic input would be the supplied project's \texttt{LFunction\_eq\_zero\_re\_of\_isPrimitive} (and RH at conductor $1$). The new theorems are not yet kernel-checked.

\paragraph{Further research questions.}
Can the near-peak return lengths be bounded explicitly in terms of $q$ and tolerance using effective simultaneous approximation, including an explicit tail bound? Can the entire support of the limiting vector be characterized without assuming independence of the ordinates? Which other Mellin kernels have real nonnegative zero weights after restricting to the critical line, and which yield stronger regularity or shorter recurrence lengths? Does a uniform version as $q$ grows lead to new arithmetic estimates? Can the deficit matrix be computed with certified error directly from finite prime-power data, yielding quantitative certificates of approximate recurrence? These are distinct from the fixed-modulus results proved here.

\begin{thebibliography}{9}
\bibitem{repo} L. Sernau et al., \emph{GRH for Dirichlet L-functions, in Lean 4}, \url{https://github.com/lsernau/grh-lean}, accessed 26 Nov. 2026. See also the file-level provenance in \texttt{SOURCES.md}.
\bibitem{wiki} Wikipedia contributors, \emph{Riemann hypothesis}, \url{https://en.wikipedia.org/wiki/Riemann_hypothesis}, accessed 26 Nov. 2026; sections on Alp\"oge's proof, explicit formulas and prime distribution.
\bibitem{davenport} H. Davenport, \emph{Multiplicative Number Theory}, 3rd ed., revised by H. L. Montgomery, Springer, 2000. Classical analytic theory of Dirichlet $L$-functions.
\bibitem{mv} H. L. Montgomery and R. C. Vaughan, \emph{Multiplicative Number Theory I: Classical Theory}, Cambridge University Press, 2007.
\bibitem{suzuki} M. Suzuki, On variants of Chebyshev's conjecture, \emph{Ramanujan J.} 68, 95 (2025), \url{https://doi.org/10.1007/s11139-025-01238-9}; preprint \url{https://arxiv.org/abs/2411.07436}. Related sign criteria with different kernels.
\bibitem{rs} M. Rubinstein and P. Sarnak, Chebyshev's bias, \emph{Experimental Mathematics} 3 (1994), 173--197, \url{https://doi.org/10.1080/10586458.1994.10504289}.
\end{thebibliography}
\end{document}
